\documentclass[a4paper]{article}
\usepackage{amsfonts}
\usepackage{amsmath}
\usepackage{amssymb}
\usepackage{graphicx}     

\newcommand{\dint}{\displaystyle\int}
\newcommand{\llim}{\lim\limits}

\begin{document}
  
\noindent \large Auteur: Sadia Nazary          \\
\noindent \large Studierichting: Informatica   \\
\noindent \large Oefeningenreeks 2B nr. 8.1    \\

\medskip

\normalsize

Bewijs dat $ \operatorname{Bgtan} \dfrac{1}{3} = \operatorname{Bgtan} \dfrac{1}{5} + \operatorname{Bgtan} \dfrac{1}{8} $\\

Rechterzijde $ = \operatorname{Bgtan}\left(\dfrac{\dfrac{1}{5} + \dfrac{1}{8}}{1 - \dfrac{1}{5} \cdot \dfrac{1}{8}}\right) $\\

$ = \operatorname{Bgtan} \left(\dfrac{\dfrac{13}{40}}{1 - \dfrac{1}{40}}\right) $\\

\medskip

$ = \operatorname{Bgtan} \left( \dfrac{13}{40} \cdot \dfrac{40}{39} \right) $\\

$ = \operatorname{Bgtan} \left( \dfrac{13}{39} \right) = \operatorname{Bgtan} \left( \dfrac{1}{3} \right) $\\

Linkerzijde $ = \operatorname{Bgtan} \left(\dfrac{1}{3}\right) $\\

$ \Rightarrow $ L.L $=$ R.L \\

Q.E.D.

\begin{thebibliography}{999}
%\bibitem{a} Referentie
\end{thebibliography}

\end{document}
